% ch5_conclusion.tex — 5장 결론
\section{결론}
\label{sec:conclusion}

\noindent{\sffamily\small\color{muted}\textbf{CHAPTER 5}}

\subsection{요약}
\label{subsec:5-1}

본 프로젝트는 ITU의 UMC(Universal Meaningful Connectivity) 프레임워크를 서울시 25개 자치구에 적용하여, 도시 내부의 디지털 격차를 측정\textperiodcentered{}분석\textperiodcentered{}해석하는 통합 분석 체계를 구축하였다. 세 단계의 분석---UMC 종합지수 구축(3.1), HLM 다수준 분석(3.2), LLM 기반 텍스트 분석(3.3)---을 순차적 삼각검증으로 결합하여, 서울의 디지털 격차가 어디에 존재하고, 왜 발생하며, 어떻게 경험되는지를 다층적으로 규명하였다.

첫째, UMC 종합지수 분석 결과, 서울 자치구 간 디지털 연결성에는 최대 2.5배의 격차가 존재하였다. 서초구(0.695)와 중랑구(0.279)의 차이는 서울이 고도로 연결된 도시라는 일반적 인식과 달리, 하위 행정단위에서는 상당한 수준의 상대적 격차가 관찰됨을 보여준다. 특히 서울 북부(도봉구, 강북구, 노원구, 중랑구)에 디지털 사막이 공간적으로 군집해 있으며, 이들은 Connectivity 차원에서는 도시 평균 이상이면서도 Devices와 Safety 차원에서 극단적 결핍을 보이는 차원별 이질성을 나타냈다.

둘째, HLM 분석은 서울의 디지털 격차가 지역 간 인프라 차이보다 개인 특성에 의해 압도적으로 구조화되어 있음을 확인하였다. ICC가 0.49\%에 불과하여 분산의 99.5\%가 개인 간 차이에 기인하였으며, 교육수준(중졸 이하 $-22.6$점)과 연령($-0.35$점/세)이 가장 강력한 예측변수였다. 그러나 교차수준 상호작용에서 공공 WiFi 밀도$\times$고령(+3.03점)과 연결성$\times$저학력(+10.94점)이 유의하게 나타나, 지역 인프라가 취약계층에게 불균형적으로 큰 혜택을 제공함이 확인되었다. 독거 고령자의 추가적 격차($-4.0$점)와 연령$\times$저학력의 가속화 패턴은 복합 취약성의 존재를 뒷받침한다.

셋째, LLM 기반 텍스트 분석은 당근 동네생활 게시글 100건에 대해 3개의 독립 추론 에이전트를 병렬 실행하여 388건의 디지털 결여 가설을 생성하였다. 정보부재(63\%), 제도적 경로 결여(61\%), 기기 접근 결여(50\%)가 주요 결여 유형으로 나타났으며, 88\%의 게시글에서 2개 이상의 결여가 중첩되었다. 이는 양적 분석이 포착하지 못하는 결여의 발현 경로---어떤 단계에서 어떤 이유로 공식 경로가 이탈되는지---를 사례 수준에서 가시화하였다.

이러한 분석 결과에 기반하여, 4장에서는 people-based(독거 고령자, 저학력 고령층, 교육 기울기 대응)와 place-based(디지털 사막 차원별 맞춤 개입, 북부 권역 연합 거버넌스)를 통합하는 정책 방향을 제시하였다. 이는 2026년 시행된 디지털포용법의 기본계획 수립 체계와 연결되며, 분석 체계 자체가 유사한 데이터 구조를 갖는 다른 도시에도 적용 가능한 절차적 일반화 가능성을 갖는다.

\subsection{한계와 향후 과제}
\label{subsec:5-2}

본 프로젝트에는 다음과 같은 한계가 존재하며, 이는 향후 연구의 방향을 제시한다.

첫째, 횡단면 데이터의 한계이다. HLM 분석에 활용된 서울서베이(2023--2024)와 UMC 지수 구축에 사용된 행정 데이터는 특정 시점의 단면을 포착한 것으로, 디지털 격차의 시간적 변화를 추적하지 못한다. 교육수준이나 연령과 같은 개인 특성이 디지털 활용에 미치는 영향이 시간에 따라 어떻게 변화하는지, 정책 개입 이후 격차가 실제로 축소되는지를 확인하려면 패널 데이터의 구축이 필요하다. 서울서베이가 매년 실시되는 행정 조사라는 점은 패널 전환의 가능성을 시사하나, 현재 표본 설계는 반복 횡단면 구조이므로 개인 수준의 추적에는 한계가 있다.

둘째, 인과 식별의 제약이다. HLM은 변수 간 관계의 방향성과 크기를 추정하지만, 관찰 데이터에 기반한 분석이므로 엄밀한 인과적 해석에는 주의가 필요하다. 예컨대, 교차수준 상호작용에서 나타난 공공 WiFi $\times$ 고령의 정적 효과가 공공 WiFi의 인과적 효과인지, 혹은 공공 WiFi 밀도가 높은 자치구의 다른 특성(복지 인프라, 인구 구성 등)이 매개하는 효과인지를 분리하기 어렵다. 향후 연구에서는 정책 변화를 외생적 충격으로 활용하는 준실험 설계(이중차분법 등)를 통해 인과 관계를 보다 엄밀하게 식별할 필요가 있다.

셋째, 텍스트 분석의 대표성 한계이다. 당근 동네생활 게시글은 플랫폼 이용자, 즉 이미 디지털 접근이 가능한 집단에 한정되므로, 디지털 소외가 가장 심각한 계층의 경험은 구조적으로 과소 대표된다. 또한 현재 분석은 100건의 게시글을 대상으로 한 Phase 1(가설 생성) 결과에 기반하며, Phase 2(판정 에이전트에 의한 삼각검증 및 최종 분류)는 아직 수행되지 않았다. 향후 전수(약 5만 건) 분석과 Phase 2 판정이 완료되면, 자치구별\textperiodcentered{}차원별로 보다 체계적인 결여 분포를 도출할 수 있을 것이다.

넷째, UMC 지수 구축의 방법론적 제약이다. 본 프로젝트에서는 6개 차원에 동일 가중치를 적용하였으나, 차원 간 상대적 중요도는 맥락에 따라 달라질 수 있다. 예컨대, 기기 보유가 보편화된 서울에서 Devices 차원의 가중치를 다른 도시와 동일하게 부여하는 것이 적절한지에 대한 검토가 필요하다. 또한 Min-Max 정규화는 극단값에 민감하므로, 자치구 수(25)가 적은 상황에서 지수의 안정성에 대한 추가 검증이 요구된다.

다섯째, 자치구 수준 분석의 공간적 해상도 한계이다. 자치구(n=25)는 서울 내 디지털 격차의 공간적 변이를 포착하기에는 다소 거친 단위이다. 동일 자치구 내에서도 행정동 간 편차가 존재할 수 있으며, 3.3절의 텍스트 분석에서 게시글에 부여된 행정동 메타데이터는 보다 세밀한 공간 분석의 가능성을 시사한다. 향후 행정동 수준의 데이터 확보가 이루어지면, 공간적 해상도를 높인 분석이 가능할 것이다.

마지막으로, 본 프로젝트의 분석 체계---UMC 프레임워크 기반 종합지수 구축, 다수준 모형을 통한 개인-지역 기여 분리, LLM 기반 질적 맥락 추출---는 서울에 특화된 처방을 넘어, 유사한 행정 데이터 구조와 디지털 플랫폼 생태계를 갖춘 다른 도시에도 적용 가능한 방법론적 프레임워크로서의 가치를 갖는다. 이 프레임워크의 외적 타당성을 검증하기 위해, 상이한 도시화 수준과 디지털 인프라 환경을 가진 지역에의 적용이 향후 과제로 남는다.
